\begin{filecontents*}[overwrite]{\jobname.ist}
headings_flag 1
numhead_positive "符号"
symhead_positive "符号"
heading_prefix "{\\bfseries "
heading_suffix "}\\nopagebreak\\par\\vskip 2pt "
delim_0 "，"
delim_1 "，"
delim_2 "，"
delim_n "，"
\end{filecontents*}


\usepackage[a4paper,margin=28mm,headheight=15pt]{geometry}
\usepackage{amsmath,amssymb,mathtools,amsthm,bm}
\usepackage{aliascnt}
\usepackage{graphicx}
\usepackage{pict2e}
\usepackage{tikz}
\usetikzlibrary{arrows.meta,calc,decorations.pathmorphing}
\usepackage{caption}
\captionsetup{labelsep=space}
\usepackage{booktabs,longtable,array}
\usepackage{multicol}
\usepackage{enumitem}
\usepackage{xcolor}
\usepackage{fancyhdr}
\usepackage{microtype}
\usepackage{imakeidx}
\makeindex[title={索引},noautomatic]
\usepackage{hyperref}
\usepackage[nameinlink,noabbrev]{cleveref}

\definecolor{stfhblue}{HTML}{1F4E79}
\hypersetup{
  unicode=true,
  colorlinks=true,
  linkcolor=stfhblue,
  citecolor=stfhblue,
  urlcolor=stfhblue,
  pdftitle={辛拓扑与弗洛尔同调：第一卷（中文私译本）},
  pdfauthor={吴龙根}
}

\setcounter{secnumdepth}{3}
\setcounter{tocdepth}{2}
\renewcommand{\contentsname}{目录}
\makeatletter
\renewcommand{\@pnumwidth}{2em}
\renewcommand{\@tocrmarg}{3em}
\makeatother
\ctexset{
  chapter = {
    name = {第,章},
    number = \arabic{chapter}
  },
  part = {
    name = {第,部分},
    number = \Roman{part}
  }
}
\renewcommand{\figurename}{图}
\renewcommand{\tablename}{表}
\setlength{\parindent}{2em}
\setlength{\parskip}{0pt}
\emergencystretch=2em

\pagestyle{fancy}
\fancyhf{}
\fancyhead[LE]{\small\thepage\quad\nouppercase{\leftmark}}
\fancyhead[RO]{\small\nouppercase{\rightmark}\quad\thepage}
\renewcommand{\headrulewidth}{0.4pt}

\newtheorem{theorem}{定理}[section]
\newaliascnt{proposition}{theorem}
\newtheorem{proposition}[proposition]{命题}
\aliascntresetthe{proposition}
\newaliascnt{lemma}{theorem}
\newtheorem{lemma}[lemma]{引理}
\aliascntresetthe{lemma}
\newaliascnt{sublemma}{theorem}
\newtheorem{sublemma}[sublemma]{子引理}
\aliascntresetthe{sublemma}
\newaliascnt{corollary}{theorem}
\newtheorem{corollary}[corollary]{推论}
\aliascntresetthe{corollary}
\newaliascnt{conjecture}{theorem}
\newtheorem{conjecture}[conjecture]{猜想}
\aliascntresetthe{conjecture}
\newaliascnt{question}{theorem}
\newtheorem{question}[question]{问题}
\aliascntresetthe{question}
\theoremstyle{definition}
\newaliascnt{definition}{theorem}
\newtheorem{definition}[definition]{定义}
\aliascntresetthe{definition}
\newaliascnt{condition}{theorem}
\newtheorem{condition}[condition]{条件}
\aliascntresetthe{condition}
\newaliascnt{example}{theorem}
\newtheorem{example}[example]{例}
\aliascntresetthe{example}
\newaliascnt{exercise}{theorem}
\newtheorem{exercise}[exercise]{习题}
\aliascntresetthe{exercise}
\newaliascnt{remark}{theorem}
\newtheorem{remark}[remark]{注}
\aliascntresetthe{remark}
\renewcommand{\proofname}{证明}

\makeatletter
\@removefromreset{equation}{section}
\makeatother
\renewcommand{\theequation}{\thesection.\arabic{equation}}

\crefname{chapter}{章}{章}
\crefname{section}{节}{节}
\crefname{theorem}{定理}{定理}
\crefname{proposition}{命题}{命题}
\crefname{lemma}{引理}{引理}
\crefname{sublemma}{子引理}{子引理}
\crefname{corollary}{推论}{推论}
\crefname{conjecture}{猜想}{猜想}
\crefname{question}{问题}{问题}
\crefname{definition}{定义}{定义}
\crefname{condition}{条件}{条件}
\crefname{example}{例}{例}
\crefname{exercise}{习题}{习题}
\crefname{remark}{注}{注}
\crefname{equation}{式}{式}
\crefname{figure}{图}{图}
\crefname{table}{表}{表}

\DeclareMathOperator{\Ham}{Ham}
\DeclareMathOperator{\Symp}{Symp}
\DeclareMathOperator{\Hameo}{Hameo}
\DeclareMathOperator{\Diff}{Diff}

\DeclareRobustCommand{\stfhIndexWidetilde}[1]{\widetilde{#1}}
\newcommand{\stfhIndexWrite}[1]{%
  \begingroup
  \let\widetilde\stfhIndexWidetilde
  \index{#1}%
  \endgroup
}
\newcommand{\stfhIndex}[3]{\stfhIndexWrite{#2@#3}}
\newcommand{\stfhSubIndex}[5]{\stfhIndexWrite{#2@#3!#4@#5}}
\newcommand{\stfhIndexRangeStart}[3]{\stfhIndexWrite{#2@#3|(}}
\newcommand{\stfhIndexRangeEnd}[3]{\stfhIndexWrite{#2@#3|)}}
\newcommand{\stfhSubIndexRangeStart}[5]{\stfhIndexWrite{#2@#3!#4@#5|(}}
\newcommand{\stfhSubIndexRangeEnd}[5]{\stfhIndexWrite{#2@#3!#4@#5|)}}

\makeatletter
\renewenvironment{theindex}{%
  \chapter*{\indexname}%
  \addcontentsline{toc}{chapter}{\indexname}%
  \markboth{\indexname}{\indexname}%
  \begingroup
  \footnotesize
  \setlength{\parindent}{0pt}%
  \setlength{\parskip}{0pt}%
  \rightskip=0pt plus 1fil\relax
  \setlength{\columnsep}{1.5em}%
  \noindent\textbf{凡例}\quad 页码均为本译本实际页码；标目按“中文（英文源词）”双语列出，纯公式符号除外。中文主标目按无声调汉语拼音逐音节排序，子标目保留所属层级。%
  \par\vspace{0.55\baselineskip}%
  \begin{multicols}{2}%
  \raggedcolumns
  \let\item\@idxitem
}{%
  \end{multicols}%
  \endgroup
}
\renewcommand{\@idxitem}{\par\hangindent=1em\hangafter=1}
\renewcommand{\subitem}{\par\hangindent=2.8em\hangafter=1\hspace*{1.4em}}
\makeatother


\newcommand{\termSymplecticManifold}{辛流形}

\newcommand{\termSymplecticForm}{辛形式}

\newcommand{\termHamiltonianDiffeomorphism}{哈密顿微分同胚}

\newcommand{\termLagrangianSubmanifold}{拉格朗日子流形}

\newcommand{\termPseudoholomorphicCurve}{伪全纯曲线}

\newcommand{\termAlmostComplexStructure}{殆复结构}

\newcommand{\termAlmostKahlerManifold}{殆凯勒流形}

\newcommand{\termFloerHomology}{弗洛尔同调}

\newcommand{\termHoferNorm}{霍费尔范数}

\newcommand{\termStableMap}{稳定映射}














\usepackage{booktabs,longtable,multicol}
\newcommand{\BookGlossEntry}[2]{%
  \par\noindent\hangindent=1em\hangafter=1 #1\quad #2\par}
\newcommand{\BookSourceIndexEntry}[2]{%
  \par\begingroup\leftskip=#1em\relax\hangindent=1em\relax
  \noindent #2\par\endgroup}
\makeatletter
\newif\ifBookIndexing
\BookIndexingfalse
\newcommand{\BookBoldPage}[1]{\textbf{\hyperpage{#1}}}
\DeclareRobustCommand{\BookRecordIndex}[1]{%
  \ifBookIndexing
    \expandafter\index\expandafter{\csname BookIndex@#1\endcsname|hyperpage}%
  \fi}
\DeclareRobustCommand{\BookIndexDefinition}[1]{%
  \ifBookIndexing
    \expandafter\index\expandafter{\csname BookIndex@#1\endcsname|BookBoldPage}%
  \fi}
\DeclareRobustCommand{\BookTerm}[2]{%
  \BookRecordIndex{#1}%
  \begingroup\let\BookRecordIndex\@gobble #2\endgroup}
\AddToHook{cmd/@outputpage/before}{\let\BookRecordIndex\@gobble}
\pdfstringdefDisableCommands{%
  \def\BookTerm#1#2{#2}%
  \let\BookRecordIndex\@gobble
  \let\BookIndexDefinition\@gobble}
\expandafter\def\csname BookIndex@T0001\endcsname{weiquanchunquxian@伪全纯曲线（pseudoholomorphic curve）}
\expandafter\let\csname BookOriginal@termPseudoholomorphicCurve\endcsname\termPseudoholomorphicCurve
\renewcommand{\termPseudoholomorphicCurve}{\BookTerm{T0001}{\csname BookOriginal@termPseudoholomorphicCurve\endcsname}}
\expandafter\def\csname BookIndex@T0002\endcsname{daifujiegou@殆复结构（almost complex structure）}
\expandafter\let\csname BookOriginal@termAlmostComplexStructure\endcsname\termAlmostComplexStructure
\renewcommand{\termAlmostComplexStructure}{\BookTerm{T0002}{\csname BookOriginal@termAlmostComplexStructure\endcsname}}
\expandafter\def\csname BookIndex@T0003\endcsname{fuluoertongdiao@弗洛尔同调（Floer homology）}
\expandafter\let\csname BookOriginal@termFloerHomology\endcsname\termFloerHomology
\renewcommand{\termFloerHomology}{\BookTerm{T0003}{\csname BookOriginal@termFloerHomology\endcsname}}
\expandafter\def\csname BookIndex@T0004\endcsname{huofeierfanshu@霍费尔范数（Hofer norm）}
\expandafter\let\csname BookOriginal@termHoferNorm\endcsname\termHoferNorm
\renewcommand{\termHoferNorm}{\BookTerm{T0004}{\csname BookOriginal@termHoferNorm\endcsname}}
\makeatother

\let\BookOriginalIndex\index
\renewcommand{\index}[1]{\ifBookIndexing\BookOriginalIndex{#1}\fi}
